% TOP_PROBLEM: {"id": 2650, "title": "Kourovka Notebook Problem 21.141", "category_id": 17, "category": {"id": 17, "name": "group_theory", "display_name": "Group Theory", "description": "Problems about groups, group actions, representations, and related algebraic structures.", "slug": "group-theory", "order_index": 17, "created_at": "2026-07-31T15:26:25.671Z"}, "status": "open", "problem_number": "KOU-21.141", "set_id": 10}
% FIRST_POSTED: 2026-10-01
% ENTRY_KIND: statement_only
% UnsolvedMath ID 2650; source catalogue code KOU-21.141
% !TeX program = lualatex
% Definition and sources only; this file is not an LLM solution attempt.
\documentclass[11pt,a4paper]{article}
\usepackage[margin=25mm]{geometry}
\usepackage{fontspec}
\setmainfont{lmroman10-regular.otf}[BoldFont=lmroman10-bold.otf,
 ItalicFont=lmroman10-italic.otf,BoldItalicFont=lmroman10-bolditalic.otf]
\usepackage{amsmath,amssymb,mathtools}
\usepackage{unicode-math}
\setmathfont{latinmodern-math.otf}
\AtBeginDocument{\let\setminus\smallsetminus}
\usepackage{xurl,hyperref}
\hypersetup{colorlinks=true,urlcolor=blue}
\setcounter{secnumdepth}{0}
\setlength{\parindent}{0pt}
\setlength{\parskip}{5pt}
\setlength{\emergencystretch}{3em}
\begin{document}
\section*{Kourovka Notebook Problem 21.141}\label{2650}
{\small UnsolvedMath ID 2650 · KOU-21.141 · Group Theory · Kourovka Notebook - New Problems, Issue 21}
\subsection{Definitions and mathematical statement}
Fix a prime $p$. A pro-$p$ group is an inverse limit of finite $p$-groups,
with its compact Hausdorff topology. Finite generation means topological
generation: the closed subgroup generated by a finite set is the whole group.
A pro-$p$ group is finitely presented if it is a quotient of a free pro-$p$
group on finitely many generators by the closed normal subgroup generated by
finitely many relators. It is coherent if every topologically finitely generated
closed subgroup is finitely presented in this sense.

Let $G_1,G_2$ be coherent pro-$p$ groups with a common closed subgroup $H$.
Assume $H$ is polycyclic in the pro-$p$ sense: it has a finite subnormal series
of closed subgroups with procyclic factors. Each factor is topologically
generated by one element, hence is finite cyclic of $p$-power order or
isomorphic to the additive group $\mathbb Z_p$.

Form the amalgamated free pro-$p$ product $G=G_1\amalg_H G_2$. This is the
pushout in pro-$p$ groups: continuous homomorphisms from $G_1,G_2$ to a
pro-$p$ group that agree on $H$ extend uniquely to a continuous homomorphism
from $G$. Equivalently, impose the identifications of the two copies of $H$
in the free pro-$p$ product using their closed normal closure.
Must $G$ be coherent? All subgroups and presentations in the conclusion
retain the topological pro-$p$ conventions above.
\subsection{Short English statement}
Is coherence preserved when two coherent pro-$p$ groups are freely amalgamated over a polycyclic pro-$p$ subgroup?
\subsection{Sources}
\begin{itemize}
\item[S1] ULAM.AI and the UnsolvedMath Contributors, supplied problems.json, record 2650 (KOU-21.141), Kourovka Notebook - New Problems, Issue 21. Catalogue identity and source transcription; CC BY 4.0.
\url{https://huggingface.co/datasets/ulamai/UnsolvedMath}.
\item[S2] The Kourovka Notebook, 21st edition, new problems, Problem 21.141. Expanded formulation of the supplied catalogue statement.
\url{https://arxiv.org/pdf/1401.0300v47}.
\end{itemize}
\end{document}
